\section{Historical mechanism evidence and controlled diagnostics}

\subsection{Historical seed-0 centroid decomposition}

The following results are historical/hypothesis-forming evidence only. They
motivate the representation/head decomposition but are not confirmatory
replications. Table~\ref{tab:supp-historical-groups} reports the retained CE
centroid strata and A/C exact-L1 outcomes.

\begin{center}
\begin{tabular}{llrrrrrr}
\toprule
Dataset & Centroid stratum & $n$ & A exact & C exact & A MAE & C MAE & $\Delta\mu_p$ \\
\midrule
RetinaMNIST & rare-end-like & 48 & 0 & 11 & 1.250 & .854 & .470 \\
RetinaMNIST & representation-inward & 18 & 0 & 0 & 1.944 & 1.500 & .338 \\
Solar & rare-end-like & 724 & 0 & 559 & 1.012 & .296 & .598 \\
Solar & representation-inward & 197 & 0 & 0 & 1.492 & 1.655 & .044 \\
\bottomrule
\end{tabular}
\captionof{table}{Historical seed-0 CE centroid strata. Centroids are
training-derived diagnostics, not causal labels.}
\label{tab:supp-historical-groups}
\end{center}

The historical original-head endpoint routes are also retained: Retina CE
seed-0 has class-4 L1 routing \([1,1,9,9,0]\) (MAE 1.700; shrinkage 1.955),
and Solar CE seed-0 has X routing \([0,23,60,838,0]\) (MAE 1.115;
shrinkage 1.136). These values describe a fixed historical population and are
not pooled with the seed-1--4 evidence.

\subsection{Historical B/D head probes}

B denotes a balanced full-head refit; D denotes a controlled-scale probe. Both
are historical mechanism controls, not candidate methods or confirmatory
comparators. Their role is limited to showing why A/C is the controlled primary
contrast: B can change directions, norms, and biases together, and D isolates
one scale manipulation. The retained Solar RPS seed-0 control includes A, B,
C, and D at fixed \(\alpha=.50\); its frozen phase note reports that scale did
not improve over direction-only C. This setting-specific historical observation
does not support a universal scale conclusion.

\subsection{Phase 3.19 factorial and Phase 3.20A severity study}

The Phase 3.19 study is one frozen Retina RPS representation (seed 0), using
training-only OOF direction adaptation. It crosses replacement-balanced versus
natural sampling with CE versus RPS adaptation objectives. Table~\ref{tab:supp-factorial}
reports the saved rare-class outcomes; it does not establish a cross-dataset
method claim.

\begin{center}
\begin{tabular}{llrrrr}
\toprule
Condition & Sampling/objective & C4 MAE & Exact & Shrinkage & $z_4-z_3$ \\
\midrule
C & balanced / CE & 1.348 & 2/66 & 1.384 & .030 \\
E & natural / CE & 1.697 & 0/66 & 1.843 & -.912 \\
F & balanced / RPS & 1.318 & 2/66 & 1.363 & .083 \\
G & natural / RPS & 1.712 & 0/66 & 1.847 & -.878 \\
\bottomrule
\end{tabular}
\captionof{table}{Phase 3.19 frozen-RPS factorial summary.}
\label{tab:supp-factorial}
\end{center}

Phase 3.20A separately retrains full CE models for supports
\(66,50,33,16,8\) across seeds 0--4; it is not numerically interchangeable
with A/C. Table~\ref{tab:supp-severity} reproduces saved five-seed means.

\begin{center}
\begin{tabular}{rrrrrr}
\toprule
$N_4$ & C4 MAE & Shrinkage & Mean $p_4$ & Severe C4 & Mean $z_4-z_3$ \\
\midrule
66 & 2.330 & 2.356 & .092 & .830 & -1.230 \\
50 & 2.280 & 2.323 & .062 & .880 & -1.615 \\
33 & 2.380 & 2.398 & .045 & .890 & -1.878 \\
16 & 2.410 & 2.472 & .025 & .970 & -2.396 \\
8 & 2.380 & 2.464 & .012 & .880 & -3.221 \\
\bottomrule
\end{tabular}
\captionof{table}{Phase 3.20A saved five-seed mean outcomes.}
\label{tab:supp-severity}
\end{center}

The controlled grid shows direct class-4 probability and adjacent-margin
degradation as support falls, while MAE, shrinkage, and severe burden are not a
clean monotonic localization response. This is a bounded full-model severity
diagnostic, not causal identification of imbalance.
